滑动窗口可用于提取时间序列的局部持续同调特征；窗口共享观测可能引入特征依赖，使随机置换所依据的联合交换性失效。本文评价“固定持续图向量表示＋最大扫描”检验的有限样本校准：先以窗口均值反例说明原始观测独立同分布不足以保证窗口序列可交换，再在固定窗口数与固定观测长度两种预算下，于 28 个零假设配置比较无约束窗口级置换（M0）与依赖乘子自助法（M1），并以均值—协方差特征（MOMENTS）作非拓扑对照。名义水平 0.05 下，两种表示的 20 个重叠配置 M0 均实质反保守（持续同调表示最高拒绝率 0.9756）；M1 在全部 56 个匹配组合中降低拒绝率，但 8 个非重叠配置均实质保守，未实现统一校准。三项预设对比（两数据生成过程等权平均）分别是重叠—非重叠 M0 拒绝率差与固定窗口数非重叠端的 M1 保守偏差，估计为 0.5354、0.5330 与 0.0266。圆到八字形的矩匹配备择下，尺寸校正功效为 1.00000 与 0.06203；整体平移负对照呈相反次序，与欧氏持续同调的平移不变性一致。降低拒绝率不等于恢复名义水平，敏感性与误报控制应分别评价。

\noindent\textbf{关键词}：持续同调；变点检测；重叠滑动窗口；交换性；有限样本校准
